Here, we focus on obsolescence to analyze how the impact of highly cited papers changes over the long term and observe their lifespan structure.
First, to determine the obsolescence state, we introduce the Obsolescence vector proposed by Sun et al.\cite{SUn1}.
The Obsolescence vector is derived from the Gini coefficient and is expressed as $O(Gs,A^-)$, but in this study, we use the form $O(Gs,A^-,A^+)$. $Gs$ corresponds to the Gini coefficient, where a smaller value indicates that citations are relatively concentrated early on. $|A^-|$ and $|A^+|$ represent segments of the area under the Lorenz curve, with their values indicating the rate of decrease and increase in citations, respectively (Figure \ref{fig:obsvec}).

